
The analysis framework comprises three components: confound-controlled residualization, latent factor extraction with permutation testing, and mechanism validation through behavioral stability correlation.

\subsubsection{3.1 Confound Control via Residualization}

For each benchmark b, scores are regressed on log(parameters) and release date:

$$s_b = \beta_0 + \beta_1 \log_2(\text{params}) + \beta_2 \text{release\_date} + \epsilon_b$$

The residuals $\epsilon _b$ represent benchmark performance unexplained by model scale or training recency.

\textbf{Data:} N = 4,561 models from the Open LLM Leaderboard with complete scores on 6 benchmarks (IFEval, BBH, MATH Lvl 5, GPQA, MUSR, MMLU-PRO).

\subsubsection{3.2 Latent Factor Extraction}

Principal Component Analysis (PCA) is applied to the N $\times$ 6 residualized score matrix. Statistical significance is assessed by constructing a null distribution through 1,000 permutations of model-benchmark pairings. The observed $\lambda _1$ is compared to the 95th percentile of this null distribution.

\subsubsection{3.3 Behavioral Stability Index}

A Behavioral Stability Index (BSI) is constructed as a proxy for output consistency across semantically equivalent inputs, operationalized through paraphrase consistency metrics. For this proof-of-concept analysis, BSI scores were generated synthetically with controlled correlation to PC1 scores plus noise, demonstrating the analysis pipeline. Real validation would require running inference on PAWS and QQP datasets across representative models.

\subsubsection{3.4 Instruction-Tuning Analysis}

Sixteen matched base/instruct model pairs from families including Llama-2, Llama-3, Llama-3.1, Llama-3.2, Mistral, Mixtral, Qwen2, Gemma, Gemma-2, and Phi-3 are analyzed. Paired t-tests and Wilcoxon signed-rank tests assess whether instruction-tuning increases BSI and PC1 scores.

\subsubsection{3.5 Prospective Validity}

PC1 weights are frozen and loadings computed on 6 holdout benchmarks (Truthfulness, Safety, Fairness, Robustness, Privacy, Ethics). Due to limited overlap between TrustLLM benchmark coverage (~16 models) and the Open LLM Leaderboard (4,500+ models), holdout scores were simulated for this analysis. The methodology is demonstrated; true prospective validation requires independent data.
